%------------------------------------------------------------------------------------
\section{Введение}

В настоящее время актуальной задачей для развития сферы телекоммуникаций
является создание группировок спутников для покрытия мобильной связью 
поверхности Земли. Такое покрытие должно существенно упростить обеспечение
связи в труднодоступных местах, тем самым обеспечив дальнейшее развитие науки и
техники.  

Современные технологии космической связи требуют высокой надежности и
эффективности передачи данных, что делает актуальным исследование
различных факторов, влияющих на качество сигналов.

Одним из таких факторов
является эффект Доплера, который возникает в результате относительного
движения источника и приемника сигналов. В современных системах передачи
данных %{\ref{steputin20176g_v2}},
использующих ортогональное частотное деление (OFDM,
Orthogonal Frequency Division Multiplexing), влияние эффекта Доплера может
существенно сказаться на характеристиках принимаемого сигнала, таких как
уровень искажений, скорость передачи данных и устойчивость к помехам.

OFDM является одной из наиболее перспективных технологий для
реализации высокоскоростной передачи информации в условиях ограниченной
полосы пропускания и сложных каналов связи.

Благодаря своей способности
эффективно противостоять межсимвольной интерференции и многолучевому
распространению, OFDM широко применяется в системах связи.

% здесь можно добавить множество ссылок на работы

Однако при
наличии значительных скоростей движения приемников или передатчиков,
таких как спутники или космические аппараты, эффект Доплера приводит к
смещению частот поднесущих, что в свою очередь вызывает проблемы с
синхронизацией и декодированием информации.

Наряду со смещением частот, для борьбы с которым есть множество методов,
возникает эффект масштабирования спектра, который приводит к возникновению
интерференции между поднесущими. Так как основной принцип OFDM состоит в ортогональности
поднесущих, данный метод особенно чувствителен к такой интерференции.

% Начинаем перечислять ссылки
Множество работ было посвящено смещению спектра, например следующие. %(\cite{maykov2014vliyanie}, \cite{panko2018poliak}, \cite{Zhuravleva2021}). 
В первой работе (Влияние эффектов Доплера на {OFDM} сигнал) было исследовано влияние смещения несущей частоты OFDM сигнала
 на сигнальное созвездие, влиянием расширения спектра в данной статье было пренебнрежено, так как относительная скорость приемника/передатчика
  была достаточно мала (по сравнению со скоростью движения спутника). Вторая работа (ИССЛЕДОВАНИЕ БИТОВЫХ ОШИБОК, ОБУСЛОВЛЕННЫХ ЭФФЕКТОМ ДОПЛЕРА)
посвящена влиянию Доплеровского сдвига в космических системах связи, было исследовано влияние частоты доплера на вероятность битовой ошибки,
а также проведено усреднение этого влияния относительно скорости спутника. Третья статья (Оценка влияния эффекта Доплера на качетво радиосвязи в условиях высокоскоростного движения)
также была посвязена оценке влияния доплеровского сдига.

% Тут можно привести описание этих работ

Следующие работы исследовали влияние масштабирования спектра и интерференции поднесущих
(
Analysis of Doppler frequency shift response mechanism based on OFDM system,
Bounds on the interchannel interference of OFDM in time-varying impairments,
ICI Performance Analysis for All Phase OFDM Systems,
Intercarrier Interference in OFDM: A deterministic model for transmissions in mobile environments with imperfect synchronization,
The effects of Doppler spreads in OFDM(A) mobile radio systems
).
%(\cite{Li2018}, \cite{YeLi2001}, \cite{Ge2009}),
%{\cite{Oberli2008}, \cite{Robertson1999}}

В данной же работе будет рассмотрено влияние масштабирования спектра
без переотражений на уровень интерференции поднесущих (ICI) и уровень битовых ошибок (BER).
Такое допущение об отсутсвии переотажений было принято в рамках системы космической связи,
когда передатчик и приемник почти всегда находятся в зоне прямой видимости (LOS), а
влиянием переотражений от окружающих обьектов можно пренебречь.

%------------------------------------------------------------------------------------
\newpage
\section{Цель}
\textbf{Цель:} В данной работе основное внимание уделяется влиянию
доплеровского масштабирования спектра в чистом виде, без учета многолучевого распространения.
Это допущение адекватно для канала "спутник-земля" в условиях прямой видимости (Line-of-Sight, LOS),
где переотражениями от окружающих объектов можно пренебречь.

\textbf{Задачи:}
\begin{enumerate}
    \item Проанализировать, как эффект масштабирования спектра, изолированный от других факторов,
     влияет на уровень интерференции между поднесущими (ICI).
    \item Оценить непосредственное воздействие данного эффекта на вероятность битовой ошибки (BER) в OFDM-системе.
    \item Дать количественную оценку деградации качества связи, вызванной исключительно высокоскоростным относительным движением в LOS-канале.
\end{enumerate}